\chapter{Những sự thật về con người}
\markboth{Những sự thật về con người}{Truths About People}


\section{Người ta thích bạn… nhưng không muốn bạn hơn họ}
\begin{truyen}{Người ta thích bạn… nhưng không muốn bạn hơn họ}{Life lesson}
\chuhoa{ộ}{t người thành đạt kể: khi anh còn làm thuê, ai cũng thân. Khi anh vượt lên, nhiều người xa dần. Họ thích anh — nhưng không muốn anh hơn họ.}
\textit{(A successful man said: when he was still an employee, everyone was close. When he rose above, many drifted away. They liked him — but didn't want him to be better than them.)}
\end{truyen}

\begin{baihoc}
Bạn thật không sợ bạn hơn họ.
\textit{(Real friends are not afraid of you being better.)}
\end{baihoc}

\begin{ghinhoanh}
\textbf{Short English to remember:}

Real friends are not afraid of you being better.

\end{ghinhoanh}

\begin{tuhocnhanh}
1. Câu chuyện này gợi cho bạn suy nghĩ gì? \textit{What does this story make you think?}

2. Bạn sẽ nhớ bài học nào nhất? \textit{Which lesson will you remember most?}
\end{tuhocnhanh}
\ngancach


\section{Khi bạn nghèo, ít người nghe bạn nói}
\begin{truyen}{Khi bạn nghèo, ít người nghe bạn nói}{Life lesson}
\chuhoa{A}{nh từng nói điều đúng nhưng không ai chịu nghe. Khi anh có địa vị, cùng câu nói ấy mọi người gật đầu. Khi nghèo, tiếng nói của anh nhỏ đi.}
\textit{(He once said the right thing but no one listened. When he had status, the same words got nods. When poor, his voice grew smaller.)}
\end{truyen}

\begin{baihoc}
Giá trị lời nói đôi khi bị gắn với vị thế — đau nhưng thật.
\textit{(The value of words is sometimes tied to status — painful but true.)}
\end{baihoc}

\begin{ghinhoanh}
\textbf{Short English to remember:}

The value of words is sometimes tied to status — painful but true...

\end{ghinhoanh}

\begin{tuhocnhanh}
1. Câu chuyện này gợi cho bạn suy nghĩ gì? \textit{What does this story make you think?}

2. Bạn sẽ nhớ bài học nào nhất? \textit{Which lesson will you remember most?}
\end{tuhocnhanh}
\ngancach


\section{Khi bạn thành công, nhiều người muốn làm bạn}
\begin{truyen}{Khi bạn thành công, nhiều người muốn làm bạn}{Life lesson}
\chuhoa{L}{úc khó khăn anh đếm được vài người. Lúc thành công, tin nhắn và lời mời không dứt. Không phải tất cả đều giả — nhưng không phải tất cả đều thật.}
\textit{(When times were hard he could count his people. When he succeeded, messages and invitations never stopped. Not all were fake — but not all were real.)}
\end{truyen}

\begin{baihoc}
Thành công kéo theo nhiều "bạn" — nhưng bạn thật thường đã có từ trước.
\textit{(Success brings many "friends" — but real friends were usually there before.)}
\end{baihoc}

\begin{ghinhoanh}
\textbf{Short English to remember:}

Success brings many "friends" — but real friends were usually the...

\end{ghinhoanh}

\begin{tuhocnhanh}
1. Câu chuyện này gợi cho bạn suy nghĩ gì? \textit{What does this story make you think?}

2. Bạn sẽ nhớ bài học nào nhất? \textit{Which lesson will you remember most?}
\end{tuhocnhanh}
\ngancach


\section{Người tốt thường bị thử thách nhiều}
\begin{truyen}{Người tốt thường bị thử thách nhiều}{Life lesson}
\chuhoa{C}{ô ấy luôn giúp người, không nói xấu ai. Vậy mà cô bị nghi ngờ, bị lợi dụng. Người tốt thường bị thử thách nhiều — không phải vì trời bất công, mà vì họ dám sống đúng.}
\textit{(She always helped, never spoke ill. Yet she was doubted and used. Good people are often tested more — not because life is unfair, but because they dare to live rightly.)}
\end{truyen}

\begin{baihoc}
Bị thử thách nhiều không có nghĩa bạn sai — đôi khi là vì bạn đúng.
\textit{(Being tested more doesn't mean you're wrong — sometimes it's because you're right.)}
\end{baihoc}

\begin{ghinhoanh}
\textbf{Short English to remember:}

Being tested more doesn't mean you're wrong — sometimes it's beca...

\end{ghinhoanh}

\begin{tuhocnhanh}
1. Câu chuyện này gợi cho bạn suy nghĩ gì? \textit{What does this story make you think?}

2. Bạn sẽ nhớ bài học nào nhất? \textit{Which lesson will you remember most?}
\end{tuhocnhanh}
\ngancach


\section{Không phải ai cười với bạn cũng thích bạn}
\begin{truyen}{Không phải ai cười với bạn cũng thích bạn}{Life lesson}
\chuhoa{A}{nh quen nhiều khuôn mặt cười chào. Đến khi anh vấp ngã, những khuôn mặt ấy biến mất. Nụ cười dễ cho — thiện chí thật thì hiếm.}
\textit{(He knew many smiling faces. When he fell, those faces vanished. A smile is easy to give — real goodwill is rare.)}
\end{truyen}

\begin{baihoc}
Cười với bạn dễ; đứng bên bạn lúc khó mới khó.
\textit{(Smiling at you is easy; standing by you when it's hard is what's hard.)}
\end{baihoc}

\begin{ghinhoanh}
\textbf{Short English to remember:}

Smiling at you is easy; standing by you when it's hard is what's ...

\end{ghinhoanh}

\begin{tuhocnhanh}
1. Câu chuyện này gợi cho bạn suy nghĩ gì? \textit{What does this story make you think?}

2. Bạn sẽ nhớ bài học nào nhất? \textit{Which lesson will you remember most?}
\end{tuhocnhanh}
\ngancach


\section{Người ta thường tin điều họ muốn tin}
\begin{truyen}{Người ta thường tin điều họ muốn tin}{Life lesson}
\chuhoa{D}{ù anh đưa ra bằng chứng, có người vẫn giữ ý cũ. Họ không tin vì sự thật — họ tin vì điều đó hợp với điều họ muốn.}
\textit{(Even with evidence, some keep their old view. They don't believe because of truth — they believe because it fits what they want.)}
\end{truyen}

\begin{baihoc}
Thay đổi niềm tin của người khác khó hơn ta tưởng.
\textit{(Changing someone's beliefs is harder than we think.)}
\end{baihoc}

\begin{ghinhoanh}
\textbf{Short English to remember:}

Changing someone's beliefs is harder than we think.

\end{ghinhoanh}

\begin{tuhocnhanh}
1. Câu chuyện này gợi cho bạn suy nghĩ gì? \textit{What does this story make you think?}

2. Bạn sẽ nhớ bài học nào nhất? \textit{Which lesson will you remember most?}
\end{tuhocnhanh}
\ngancach


\section{Lời nói ngọt đôi khi không thật}
\begin{truyen}{Lời nói ngọt đôi khi không thật}{Life lesson}
\chuhoa{N}{hững lời khen dễ nghe, những lời hứa đẹp — anh từng nhận nhiều. Khi cần giữ lời, nhiều người im. Lời ngọt dễ nói; việc làm mới thật.}
\textit{(Sweet words and pretty promises — he had received many. When it was time to keep their word, many fell silent. Sweet words are easy; actions are what's real.)}
\end{truyen}

\begin{baihoc}
Đừng tin quá nhanh vào lời ngọt — hãy xem việc làm.
\textit{(Don't trust sweet words too quickly — watch the actions.)}
\end{baihoc}

\begin{ghinhoanh}
\textbf{Short English to remember:}

Don't trust sweet words too quickly — watch the actions.

\end{ghinhoanh}

\begin{tuhocnhanh}
1. Câu chuyện này gợi cho bạn suy nghĩ gì? \textit{What does this story make you think?}

2. Bạn sẽ nhớ bài học nào nhất? \textit{Which lesson will you remember most?}
\end{tuhocnhanh}
\ngancach


\section{Người im lặng thường hiểu nhiều hơn}
\begin{truyen}{Người im lặng thường hiểu nhiều hơn}{Life lesson}
\chuhoa{T}{rong phòng họp, người nói nhiều nhất không phải người hiểu nhất. Người ngồi im lắng nghe, khi lên tiếng thường đúng trọng tâm. Im lặng cho phép họ nghe và nghĩ.}
\textit{(In the meeting, the one who spoke most was not the one who understood most. The one who sat silent and listened often hit the point when they spoke. Silence let them hear and think.)}
\end{truyen}

\begin{baihoc}
Im lặng không phải thiếu ý kiến — đôi khi là đang tích lũy sự hiểu.
\textit{(Silence isn't lack of opinion — sometimes it's gathering understanding.)}
\end{baihoc}

\begin{ghinhoanh}
\textbf{Short English to remember:}

Silence isn't lack of opinion — sometimes it's gathering understa...

\end{ghinhoanh}

\begin{tuhocnhanh}
1. Câu chuyện này gợi cho bạn suy nghĩ gì? \textit{What does this story make you think?}

2. Bạn sẽ nhớ bài học nào nhất? \textit{Which lesson will you remember most?}
\end{tuhocnhanh}
\ngancach


\section{Người giúp bạn lúc khó khăn mới là bạn thật}
\begin{truyen}{Người giúp bạn lúc khó khăn mới là bạn thật}{Life lesson}
\chuhoa{K}{hi anh còn nghèo, rất ít người hỏi thăm. Khi anh thành công, nhiều người muốn làm bạn. Nhưng người đã giúp anh khi anh khó khăn — anh nhớ suốt đời.}
\textit{(When he was still poor, very few asked after him. When he succeeded, many wanted to be friends. But the one who helped him when he was in need — he remembers for life.)}
\end{truyen}

\begin{baihoc}
Bạn thật thường xuất hiện lúc mình yếu nhất.
\textit{(Real friends often appear when you're at your weakest.)}
\end{baihoc}

\begin{ghinhoanh}
\textbf{Short English to remember:}

Real friends often appear when you're at your weakest.

\end{ghinhoanh}

\begin{tuhocnhanh}
1. Câu chuyện này gợi cho bạn suy nghĩ gì? \textit{What does this story make you think?}

2. Bạn sẽ nhớ bài học nào nhất? \textit{Which lesson will you remember most?}
\end{tuhocnhanh}
\ngancach


\section{Ai cũng có nỗi lo riêng}
\begin{truyen}{Ai cũng có nỗi lo riêng}{Life lesson}
\chuhoa{A}{nh từng nghĩ chỉ mình anh vất vả. Sau này anh mới biết: người bên cạnh cũng có nỗi lo, nỗi sợ, nỗi đau. Ai cũng đang gồng gánh một phần đời.}
\textit{(He once thought only he had it hard. Later he learned: the person next to him had worries, fears, and pain too. Everyone is carrying part of life.)}
\end{truyen}

\begin{baihoc}
Đừng nghĩ chỉ mình mình khó — mở mắt ra sẽ thấy ai cũng có gánh.
\textit{(Don't think only you have it hard — open your eyes and you'll see everyone carries a load.)}
\end{baihoc}

\begin{ghinhoanh}
\textbf{Short English to remember:}

Don't think only you have it hard — open your eyes and you'll see...

\end{ghinhoanh}

\begin{tuhocnhanh}
1. Câu chuyện này gợi cho bạn suy nghĩ gì? \textit{What does this story make you think?}

2. Bạn sẽ nhớ bài học nào nhất? \textit{Which lesson will you remember most?}
\end{tuhocnhanh}
\ngancach
